% ============================================
% Assistant/Associate Professor Cover Letter – Carnegie Mellon University (Tepper School of Business)
% ============================================

\documentclass[11pt,letterpaper]{letter}

\usepackage{geometry}
\usepackage{microtype}
\usepackage[hidelinks]{hyperref}

\geometry{left=25mm,right=25mm,top=25mm,bottom=25mm}

\signature{Decory Edwards}
\address{
Department of Economics \\
Johns Hopkins University \\
Baltimore, MD 21218 \\
\texttt{dedwar65@jh.edu}
}

\begin{document}

\begin{letter}{Search Committee \\
Economics Faculty Search \\
Tepper School of Business \\
Carnegie Mellon University \\
Pittsburgh, PA 15213 \\ USA}

\opening{Dear Members of the Search Committee,}

I am writing to apply for the tenure-track faculty position in Economics in the Tepper School of Business at Carnegie Mellon University. I am a Ph.D. candidate in Economics at Johns Hopkins University, advised by Professors Christopher D.~Carroll, Nicholas W.~Papageorge, and Stelios Fourakis, and I anticipate completing my degree in May 2026. My work lies at the intersection of macroeconomics, wealth and income dynamics, and heterogeneous-agent modeling, with an emphasis on drawing empirical insights from micro data to inform broader macroeconomic behavior.

My research agenda centers on understanding how heterogeneity in returns to wealth shapes aggregate outcomes and distributional dynamics. In my job-market paper, I estimate persistent heterogeneity in household-level portfolio returns using data from the Survey of Consumer Finances and embed those estimates in a structural heterogeneous-agent model. I find that accounting for return heterogeneity markedly improves the model’s ability to match empirical wealth distributions, revealing key channels through which risk and return disparities impact macroeconomic inequality. In a second project, I analyze the Health and Retirement Study to uncover a hump-shaped relationship between interpersonal trust and realized portfolio returns, suggesting a behavioral pathway that informs both saving and investment decisions as well as exposure to aggregate risk.

Carnegie Mellon’s Tepper School of Business, with its distinctive blend of economics, data science, and quantitative rigor, is an ideal environment for this research agenda. The opportunity to contribute to a department that spans undergraduate, MBA, and Ph.D. instruction aligns closely with my intellectual interests and teaching background. I am particularly drawn to Tepper’s emphasis on analytical approaches to economic questions, its interdisciplinary strengths, and its tradition of integrating cutting-edge research with rigorous pedagogy.

\textbf{Teaching.} My teaching experience includes leading discussion sections and designing problem sets for undergraduate \textit{Principles of Macroeconomics}, \textit{Intermediate Macroeconomic Theory}, and an upper-level \textit{Macroeconomic Policy} seminar. My approach emphasizes clarity in exposition, structured problem solving, and helping students transfer theoretical insights to empirical analysis. At Tepper, I am prepared to teach core undergraduate economics courses as well as graduate-level theory and applied courses that draw on modern macroeconomic and quantitative methods. I welcome the opportunity to mentor students across all programs and to support Tepper’s mission of fostering analytical excellence and real-world problem solving.

I believe I would be a strong fit for Carnegie Mellon University because my research agenda engages deeply with quantitative modeling and empirical methods that complement the school’s strengths, and because I am committed to delivering teaching that integrates rigorous analytical foundations with applications that resonate with students across business and economics curricula. I would be honored to contribute to the Tepper School’s ongoing tradition of innovation in research and education.

Thank you for your consideration. I welcome the opportunity to discuss my application and qualifications further.

\closing{Sincerely,}

\end{letter}
\end{document}
